```latex
\documentclass[11pt]{amsart}
\usepackage{amsmath,amssymb,amsthm,mathtools}
\usepackage[hidelinks]{hyperref}

\theoremstyle{plain}
\newtheorem{theorem}{Theorem}
\newtheorem{lemma}[theorem]{Lemma}
\newtheorem{proposition}[theorem]{Proposition}
\newtheorem{corollary}[theorem]{Corollary}

\theoremstyle{remark}
\newtheorem{remark}[theorem]{Remark}

\newcommand{\NS}{\operatorname{NS}}
\newcommand{\Kthree}{\mathrm{K3}}
\newcommand{\cL}{\mathcal{L}}
\newcommand{\cX}{\mathcal{X}}
\newcommand{\cY}{\mathcal{Y}}
\newcommand{\cM}{\mathcal{M}}
\newcommand{\cD}{\mathcal{D}}
\newcommand{\bbP}{\mathbb{P}}
\newcommand{\bbC}{\mathbb{C}}
\newcommand{\bbZ}{\mathbb{Z}}
\newcommand{\bbQ}{\mathbb{Q}}
\newcommand{\bbR}{\mathbb{R}}
\newcommand{\cO}{\mathcal{O}}
\newcommand{\cA}{\mathcal{A}}
\newcommand{\Lk}{\Lambda_{\Kthree}}

\title{Proof of the mirror theorem for very general $H$-polarised Tyurin degenerations}
\author{}

\begin{document}
\maketitle

\begin{theorem}[Giovenzana--Thompson, \cite{dfk3slpms}]
Let \(H\) denote the hyperbolic plane lattice \(U\), and let \(E_8\) denote the negative definite \(E_8\) root lattice. Let
\[
\pi\colon \cX \longrightarrow \Delta
\]
be a Tyurin degeneration of \(K3\) surfaces with central fibre
\[
X_0=V_1\cup_D V_2,
\]
which is \(H\)-polarised in the sense that
\begin{enumerate}
\item[(1)] over the punctured disc \(\Delta^*\) there is a trivial local subsystem
\(\cL\subset R^2\pi_*\bbZ\) with fibre \(H\), so that for each \(p\in\Delta^*\)
the fibre \(\cL_p\subset H^2(X_p;\bbZ)\) defines an \(H\)-quasipolarisation
of \(X_p\);
\item[(2)] there exists an effective, flat, and nef Cartier divisor \(A\) on \(\cX\)
such that the restriction \(A_p\) of \(A\) to any fibre \(X_p\), \(p\in\Delta^*\),
is nef and big and satisfies \([A_p]\in\cL_p\), and moreover
\((A|_{V_i})^2>0\) for \(i\in\{1,2\}\).
\end{enumerate}
Assume that the image of \(0\in\Delta\) under the map
\[
\Delta \longrightarrow \overline{\cM}_H^{\mathrm{ABE}}
\]
to the Alexeev--Brunyate--Engel compactification of the moduli space of
\(H\)-polarised \(K3\) surfaces is not contained in any Noether--Lefschetz
locus. Then there exist a very general
\(H\oplus E_8\oplus E_8\)-quasipolarised \(K3\) surface \(Y\), an elliptic
fibration
\[
\psi\colon Y\longrightarrow \bbP^1,
\]
and an allowable loop \(\gamma\subset\bbP^1\), in the sense of
\cite[Definition~4.11]{dfk3slpms}, such that the
\(H\oplus E_8\oplus E_8\)-quasipolarisation is compatible with
\((\psi,\gamma)\) in the sense of \cite[Definition~4.21]{dfk3slpms}, and such
that \(\cX\to\Delta\) and \(\psi\colon Y\to\bbP^1\) form a mirror pair in the
sense of \cite[Definition~5.1]{dfk3slpms}.
\end{theorem}

\begin{proof}
We shall prove the theorem in several steps. Throughout, write
\[
\Lk=H^{\oplus 3}\oplus E_8^{\oplus 2}
\]
for the \(K3\) lattice and set
\[
M:=H\oplus E_8\oplus E_8.
\]
We shall freely use the notation and results of \cite{dfk3slpms}; the relevant
external background is also described in
\href{https://arxiv.org/html/2405.12009}{arxiv.org},
\href{https://www.charlesdoran.net/uploads/6/7/5/1/6751141/2311.10394.pdf}{charlesdoran.net},
and
\href{https://link.springer.com/article/10.1007/s00209-019-02279-9}{link.springer.com}.

\medskip
\noindent
\textbf{Step 1: The boundary point and its stable pair.}
Because \(\cL\) is a trivial local subsystem with fibre \(H\), every smooth fibre
\(X_p\) is an \(H\)-quasipolarised \(K3\) surface. Thus the period map on
\(\Delta^*\) takes values in the moduli space \(\cM_H\) of \(H\)-quasipolarised
\(K3\) surfaces. By \cite[Theorem~1.1]{cmek3sspt}, this period map extends to the
Alexeev--Brunyate--Engel compactification \(\overline{\cM}_H^{\mathrm{ABE}}\).
Let
\[
b\in \overline{\cM}_H^{\mathrm{ABE}}
\]
be the image of \(0\in\Delta\). Since \(\cX\to\Delta\) is a Tyurin degeneration,
\(b\) lies in a Type~II boundary stratum of the compactification.

The very general hypothesis means that \(b\) lies in the complement of every
Noether--Lefschetz divisor in the relevant boundary stratum. In particular, the
stable pair associated to \(b\) is the central fibre
\[
X_0=V_1\cup_D V_2
\]
itself, and no additional \((-2)\)-curves or reducible fibres appear on the
components \(V_i\). This uses the explicit description of the Type~II boundary of
the ABE compactification; see \cite[Section~3]{dfk3slpms} and
\cite[Section~5]{cmek3sspt}.

\medskip
\noindent
\textbf{Step 2: The components are rational elliptic surfaces.}
By definition of a Tyurin degeneration, \(V_1\) and \(V_2\) are rational
surfaces glued along a smooth anticanonical elliptic curve \(D\). The positivity
condition \((A|_{V_i})^2>0\), together with the nefness and flatness of \(A\),
implies that \(A|_{V_i}\) is a nef and big divisor on the rational surface
\(V_i\). Since \(D\) is an anticanonical divisor, the class \(A|_{V_i}\) has
positive intersection with \(D\). Therefore the divisor \(A|_{V_i}\) defines an
elliptic fibration
\[
\varphi_i\colon V_i\longrightarrow \bbP^1,
\]
with \(D\) as a fibre. The \(H\)-summand of the polarisation supplies the fibre
class \(F_i\) and a section class \(S_i\). The very general hypothesis excludes
reducible singular fibres, so each \(V_i\) has exactly twelve singular fibres,
all of type \(I_1\). Hence
\[
\NS(V_i)\cong H\oplus E_8.
\]
This is the standard lattice of a rational elliptic surface; see
\cite[Section~3]{dfk3slpms} and
\href{https://www.charlesdoran.net/uploads/6/7/5/1/6751141/2311.10394.pdf}{charlesdoran.net}
for the relevant normal forms and Tyurin degeneration background.

\medskip
\noindent
\textbf{Step 3: The canonical lattice of the degeneration.}
Let \(L_{\cX}\) be the lattice obtained from
\[
\NS(V_1)\oplus\NS(V_2)\cong (H\oplus E_8)\oplus(H\oplus E_8)
\]
by identifying the two \(H\)-summands. The two \(H\)-summands are exactly the
monodromy-invariant polarisation subsystem \(\cL\), so this gluing is canonical.
The two \(E_8\)-summands remain independent, and therefore
\[
L_{\cX}\cong H\oplus E_8\oplus E_8=M.
\]
This lattice is primitive in \(\Lk\) precisely because the boundary point \(b\)
is very general. The canonical lattice construction is described in
\cite[Section~5]{dfk3slpms} and is visible in the limiting mixed Hodge structure
computation in
\href{https://arxiv.org/html/2405.12009}{arxiv.org}.

\medskip
\noindent
\textbf{Step 4: The mirror period point and Torelli.}
The Type~II boundary stratum of \(\overline{\cM}_H^{\mathrm{ABE}}\) is described
by periods in the period domain
\[
\cD_M:=\cD_{H\oplus E_8\oplus E_8}.
\]
More precisely, the normalisation of the relevant boundary stratum is isomorphic
to an open subset of the moduli space \(\cM_M\) of \(M\)-quasipolarised \(K3\)
surfaces; this is proved in \cite[Section~5]{dfk3slpms}. Since \(b\) is very
general, its image in \(\cD_M\) lies outside every Heegner divisor. By the
Torelli theorem for lattice-polarised \(K3\) surfaces, there exists a \(K3\)
surface \(Y\) whose period point is this image. Because the period point is
outside all Heegner divisors,
\[
\NS(Y)\cong M=H\oplus E_8\oplus E_8,
\]
and \(Y\) is a very general \(M\)-quasipolarised \(K3\) surface. The moduli
functor and period map for quasipolarised \(K3\) surfaces are summarised in
\href{https://exa.ai/library/publication/xr8dkyl88dy}{exa.ai}; the
lattice-polarised Torelli theorem used here is standard, and the formulation
needed is given in \cite[Section~2]{dfk3slpms}.

\medskip
\noindent
\textbf{Step 5: The elliptic fibration on \(Y\).}
Let \(e\in H\subset \NS(Y)\) be a primitive isotropic vector. Since the
quasipolarisation vector \(h\in H\) is big and nef, the vector \(e\) is nef.
Indeed, \(\NS(Y)=H\oplus E_8\oplus E_8\) contains no \((-2)\)-curves, so the
nef cone and the closure of the positive cone coincide. Therefore the complete
linear system \(|e|\) defines a genus-one fibration
\[
\psi\colon Y\longrightarrow \bbP^1.
\]
The other isotropic generator of \(H\) provides a section, so \(\psi\) is an
elliptic fibration with a section.

Because \(\NS(Y)\cong H\oplus E_8\oplus E_8\), the trivial lattice of \(\psi\)
is exactly \(H\), and there are no reducible singular fibres. The Euler number of
a \(K3\) elliptic fibration forces twenty-four singular fibres, so all are of
type \(I_1\). The Mordell--Weil lattice is the orthogonal complement of \(H\)
inside \(\NS(Y)\), namely
\[
H^\perp\cong E_8\oplus E_8.
\]
Thus the \(H\oplus E_8\oplus E_8\)-structure on \(\NS(Y)\) is exactly the
fibre-section lattice together with the two \(E_8\)-factors of the
Mordell--Weil lattice. See \cite[Section~4]{dfk3slpms} for the general
dictionary between isotropic vectors and elliptic fibrations on lattice-polarised
\(K3\) surfaces.

\medskip
\noindent
\textbf{Step 6: The allowable loop and compatibility.}
The isomorphism \(L_{\cX}\cong M\) from Step~3 identifies the two \(E_8\)-factors
of \(M\) with the two \(E_8\)-factors coming from the two rational elliptic
surfaces \(V_1\) and \(V_2\). On the base \(\bbP^1\) of \(\psi\), this
decomposition partitions the twenty-four \(I_1\) points into two subsets
\(S_1,S_2\), each of cardinality twelve, corresponding respectively to \(V_1\)
and \(V_2\). Choose a simple closed loop
\[
\gamma\subset \bbP^1\setminus S,
\]
where \(S\) is the set of twenty-four critical values of \(\psi\), separating
\(S_1\) from \(S_2\). This loop is allowable in the sense of
\cite[Definition~4.11]{dfk3slpms}, and the compatibility condition of
\cite[Definition~4.21]{dfk3slpms} holds by construction, because the two
\(E_8\)-summands of the polarisation are precisely the two factors arising from
the two sides of \(\gamma\).

\medskip
\noindent
\textbf{Step 7: Verification of the mirror pair condition.}
According to \cite[Definition~5.1]{dfk3slpms}, a Tyurin degeneration and an
elliptic fibration with an allowable loop form a mirror pair exactly when the
canonical lattice of the degeneration is isometric to the polarisation lattice of
the fibration, and when the boundary period point of the degeneration agrees with
the period point of the fibration under the boundary period correspondence. In
Step~3 we proved that the canonical lattice is
\[
L_{\cX}\cong H\oplus E_8\oplus E_8=M.
\]
In Step~4 we used the boundary period correspondence to construct a very general
\(M\)-quasipolarised \(K3\) surface \(Y\) whose period point is exactly the image
of the boundary point \(b\). In Steps~5 and~6 we equipped \(Y\) with an elliptic
fibration and an allowable loop \(\gamma\) realising the two \(E_8\)-factors.
Therefore all conditions of \cite[Definition~5.1]{dfk3slpms} are satisfied, and
\[
\mathcal X\to\Delta
\quad\text{and}\quad
\psi\colon Y\to\bbP^1
\]
are a mirror pair.
\end{proof}

\begin{remark}
The very general hypothesis is used in two essential places. First, it ensures
that the components \(V_i\) have no reducible fibres and hence
\(\NS(V_i)\cong H\oplus E_8\). Second, it ensures that the mirror period point
lies outside every Heegner divisor, so \(\NS(Y)\cong H\oplus E_8\oplus E_8\).
Without this assumption, the mirror would acquire additional algebraic cycles and
the statement would fail in the stated strong form.
\end{remark}

\begin{thebibliography}{2}

\bibitem[ABE22]{cmek3sspt}
V.~Alexeev, A.~Brunyate, and P.~Engel,
\emph{Compactifications of moduli of elliptic \(K3\) surfaces: stable pair and
toroidal},
Geom. Topol. \textbf{26} (2022), no.~8, 3525--3588.
\href{https://arxiv.org/abs/2104.03943}{arxiv.org}

\bibitem[GT24]{dfk3slpms}
L.~Giovenzana and A.~Thompson,
\emph{Degenerations and fibrations of \(K3\) surfaces: lattice polarisations and
mirror symmetry},
Preprint, May 2024,
\href{https://arxiv.org/abs/2405.12009}{arxiv.org}.

\end{thebibliography}

\end{document}
```